\documentclass[10pt,a4paper]{amsart}
\usepackage[margin=19mm]{geometry}
\usepackage{amsmath,amssymb,mathtools}
\usepackage[scr=rsfs,cal=euler]{mathalfa}
\usepackage{xfrac}
\usepackage[unicode,hidelinks,bookmarks=false]{hyperref}
\newcommand{\ie}{i.e.\ }
\newcommand{\cf}{cf.\ }
\newcommand{\resp}{respectively\ }
\newcommand{\Subsection}{\subsection}
\providecommand{\nobreakdash}{\nobreak\hbox{-}}
\setlength{\emergencystretch}{2em}
\multlinegap0pt
\usepackage{tikz}
\usetikzlibrary{cd}
\usepackage{tikz}
\usepackage{amssymb}
\usepackage{mathtools}
\newtheorem{corollary}{Corollary}
\newtheorem{theorem}{Theorem}
\newtheorem{lemma}{Lemma}
\newtheorem{proposition}{Proposition}
\newcommand{\Set}{\mathbf{Set}}
\newcommand{\funct}[1]{\operatorname{\mathtt{#1}}}
\newcommand{\pset}{\mathbf{set}^\mathrm{pt}}
\DeclareMathOperator{\rep}{rep}
\newcommand{\set}{\mathbf{set}}
\newcommand{\vectF}{\mathbf{vect}_F}
\title{On the additive image of 0th persistent homology}
\author{Ulrich Bauer, Magnus Bakke Botnan, Steffen Oppermann, Johan Steen}
\date{Source: arXiv:2501.09132v3 (CC BY 4.0)}
\begin{document}
\maketitle

\noindent\emph{Source and licence.} On the additive image of 0th persistent homology. Version arXiv:2501.09132v3 (CC BY 4.0), distributed by arXiv under the Creative Commons Attribution 4.0 International licence, http://creativecommons.org/licenses/by/4.0/, as recorded in the arXiv licence metadata for this exact version and shown on the abstract page https://arxiv.org/abs/2501.09132v3. Full licence notice: \texttt{licenses/arxiv-2501.09132-CC-BY-4.0.txt}.\par\smallskip

\noindent\textit{Editorial scope.} Counted unit: the article's theorem thm.onlythin (source0.tex lines 1232--1239). The article's complete original proof of it is delivered with it (source0.tex lines 1276--1314, one \texttt{proof} environment of 39 lines, which the article places away from the statement). No mathematics is written, repaired or re-derived: the statement and the proof are the article's own lines, in the article's own notation, and no retained line is substituted or re-flowed.  Retained uncounted as the source-local context the proof needs: lines 915--922; lines 1265--1267; lines 1347--1361; lines 1452--1465.  Nothing was omitted from any retained span, so the package carries no omission record.  No retained line contains a citation, so the wrapper carries no bibliography and no bibliography entry.  The retained lines contain the article's own diagram code, which the wrapper typesets with the standard tikz packages; no image file is delivered and the package declares no asset dependency.  Editorial formatting only, in addition: an AMS \texttt{amsart} preamble that restates the article's own declarations and macros used by the retained lines, namely the environments \texttt{corollary}, \texttt{theorem}, \texttt{lemma}, \texttt{proposition} on separate counters, together with the standard packages those lines need.  The retained environments print in this document as Corollary 1, Theorem 1, Lemma 1, Proposition 1 in order of appearance, whereas the article prints the same items with the numbering of its own full document.  The complete native source of the article as distributed in the arXiv e-print for this exact version is bundled unchanged as \texttt{sources/171651/source0.tex}.
\begin{corollary} \label{cor.subcategory}
Let \( S \) be a full subcategory of \( X \).
\begin{itemize}
\item If all indecomposable additively \( \Set \)-realizable representations of \( X \) are indicator representations, then the same is true for \( S \).
\item If there are only finitely many indecomposable additively \( \Set \)-realizable representations of \( X \), then the same is true for \( S \).
\item If all linear representations of \( X \) are additively \( \Set \)-realizable, then the same is true for \( S \).
\end{itemize}
\end{corollary}

\begin{lemma} \label{lemma:indicator_for_monoids}
Let \( X \) be a finite category with one object. Then all indecomposable additively \( \Set \)-realizable representations are indicator representations only if the category has no morphisms apart from the identity.
\end{lemma}

\begin{proposition} \label{prop:when_is_Dn_realizable}
For a quiver \( Q \) of type \( D_n \), the following are equivalent.
\begin{enumerate}
\item the functor \( \funct{free}^\mathrm{pt} \colon \rep( Q, \pset) \to \rep( Q, \vectF) \)
is essentially surjective;
\item the functor \( \funct{free} \colon \rep( Q, \set) \to \rep( Q, \vectF) \)
is additively surjective;
\item not every arm of the quiver contains an ``ingoing arrow'', that is \( Q \) is not of the form
\[ \begin{tikzcd}[sep=5mm]
\circ \ar[r,-,dotted,very thick] & \circ \ar[r] & \circ \ar[r,-,dotted,very thick] & \circ & \circ \ar[l] \\
&&& \circ \ar[u] 
\end{tikzcd} \]
where the arrows in the dotted part may have any orientation.
\end{enumerate}
\end{proposition}

\begin{theorem}\label{thm.Ainf}
Let \( Q \) be a quiver of type \( \tilde A \). An indecomposable \( R \in \rep( Q, \vectF) \) is in the additive image of the functor \( \funct{free}_* \colon \rep( Q, \set ) \to \rep( Q, \vectF ) \) if and only if
\begin{itemize}
\item $R$ is a string representation, or
\item $R$ is a band representation with $V\cong F[X]/(f(X)^m)$ where
\begin{enumerate}
\item $f(X)$ is the minimal polynomial of a root of unity,
\item the number $m$ satisfies
\[ m = \begin{cases} 1 & \text{if ${\rm char}(F) = 0$} \\ {\rm char}(F)^k & \text{if ${\rm char}(F) > 0$}\end{cases},\]
where $k$ is a non-negative integer and ${\rm char}(F)$ denotes the characteristic of $F$, 
\item the isomorphism $\phi\colon V\to V$ in the definition of a band representation is given by $\phi(q) = X\cdot q$.
\end{enumerate}
\end{itemize}
\end{theorem}

\begin{theorem}
\label{thm.onlythin}
Let \( X \) be a finite category. Then every indecomposable additively \( \Set \)-realizable representation is an indicator representation if and only if \( X \) is equivalent to the poset category of a poset that is a disjoint union of
\begin{enumerate}
\item posets with a Hasse diagram of type \( A \),
\item posets with a single maximal element, such that the poset without the maximal element is a disjoint union of posets with a Hasse diagram of type \( A \) with unique minima.
\end{enumerate}
\end{theorem}

\begin{proof}[Proof of \( \Rightarrow \) in Theorem~\ref{thm.onlythin}]
By Lemma~\ref{lemma:indicator_for_monoids}, in conjunction with Corollary~\ref{cor.subcategory},  we can assume that the only morphisms from an object of \( X \) to itself are identities. If there are two objects with morphism between them in both directions then they thus necessarily compose to identities, so the two objects are isomorphic and up to equivalence we can remove one of them. Thus, up to equivalence we can assume \( X \) to not have any cycles of (non-identity) morphisms.

Next we observe that the category given by the Kronecker quiver $\begin{tikzcd} \circ\ar[r, bend left]\ar[r, bend right] & \circ \end{tikzcd}$ has non-indicator additively \( \Set \)-realizable representations, e.g.,

\[
\begin{tikzcd}[sep=7mm,ampersand replacement=\&]
  \{a,b\} \& \& \& \& F^2 \\
 ~\ar[rrrr,mapsto, shorten <= 30pt, shorten >= 30pt, "\funct{free}_*"] \& \&  \& \& ~\\
  \{a\} \arrow[uu,swap, bend right, "a\mapsto b"]\arrow[uu,"a\mapsto a", bend left]  \& \& \&  \& F\arrow[uu, "{\begin{bmatrix}0 \\ 1\end{bmatrix}}", bend left] \arrow[uu, swap, "{\begin{bmatrix}1 \\ 0\end{bmatrix}}", bend right]
\end{tikzcd}
\]
In fact, all preprojective representations of the Kronecker quiver are realizable.
This example can easily be extended to any number of parallel arrows, and using Corollary~\ref{cor.subcategory} also to any other category with at least two parallel arrows between some pair of distinct objects.
It now follows from Corollary~\ref{cor.subcategory} that there is at most one morphism between any two objects in \( X \), that is, \( X \) is a poset category.

Next, it is easy to check that any quiver of type \( D_4 \), except the one with all arrows pointing towards the middle, has a non-indicator  additively \( \Set \)-realizable representation; see  Proposition~\ref{prop:when_is_Dn_realizable}. It follows that
\begin{itemize}
\item No \( x  \in X \) is smaller than \( 3 \) or more pairwise incomparable elements of \( X \).
\item If \( x \in X \) is smaller than \( 2 \) incomparable elements then \( x \) is minimal.
\item If \( x \in X \) is greater than \( 2 \) incomparable elements then \( x \) is maximal.
\end{itemize}
Finally note that if an element is smaller than \( 2 \) incomparable elements, and one of those two in turn is bigger than at least \( 3 \) incomparable elements, then \( X \) has a subposet
\[ \begin{tikzcd}[column sep=5mm,row sep=5mm]
\circ && \circ \\
& \circ \ar[lu] \ar[ru] & \circ \ar[u] & \circ \ar[lu]
\end{tikzcd} \]
This poset has non-indicator additively \( \Set \)-realizable representations, e.g.,
\[ \begin{tikzcd}[column sep=8mm,row sep=10mm,ampersand replacement=\&]
F \& F^2  \ar[l, "\left({\begin{smallmatrix}1 & 1\end{smallmatrix}}\right)"] \ar[r, "\text{id}"]\&  F^2 \\
\& \& F \ar[u,"\left({\begin{smallmatrix}1 \\ 0 \end{smallmatrix}}\right)"] \& F \ar[lu,swap,"\left({\begin{smallmatrix}0 \\ 1\end{smallmatrix}}\right)"]
\end{tikzcd} \]

Therefore this cannot happen under the assumption that \( X \) does not have non-indicator additively \( \Set \)-realizable representations.

It follows that if there is an element in \( X \) which is smaller than two distinct maxima then \( X \) is of type \( A \) or \( \widetilde{A} \). Considering preprojective representations one easily sees that algebras of type \( \widetilde{A} \) have non-indicator \( \Set \)-realizable representations; e.g., the representation of the Kronecker quiver above easily extends to $\widetilde{A}$; see also Theorem~\ref{thm.Ainf}.

If \( X \) does not have an element which is smaller than two distinct maximal elements, and is connected, then \( X \) has a unique maximal element. Again invoking Corollary~\ref{cor.subcategory} we know that the same results hold for \( X' = X \setminus \{ \text{maximal element} \} \), with the additional condition that the elements smaller than any maximal element of \( X' \) are linearly ordered by the third bullet point above. By the first bullet point, at most two chains of descending elements from distinct maximal elements in $X'$ can intersect, and by the second bullet point, the point of intersection must be at a minimal element. It follows that \( X' \) is a disjoint union of type \( A \) quivers with unique sources.
\end{proof}

\end{document}
